A linear probe reads a clinical concept from a vision-language model's visual representation at the visual block its language model consumes. This readability promises a handle for selectively changing the model's answers by writing the probe direction back into that representation. Evaluations measure whether models encode clinical findings and answer questions about them, leaving the selectivity of such writes untested. We test readability, answerability and selective control across \cfCheckpoints{} checkpoints on NIH ChestX-ray14, CheXpert Plus and COCO using fixed-dose writes, random and sham references, and competing concept directions. Clinical concepts are readable in \cfChestReadablePct{}\% of chest checkpoint--concept cells and answerable in \cfChestAnswerablePct{}\%, but owned in \cfChestOwnedPct{}\%. At the fixed dose, the probe direction clears random and sham references in \cfChestRefMet{} of the chest cells, and \cfChestOwned{} of those are owned; most chest writes therefore fail the reference, and a further group that clears it loses to a competing clinical direction. Readability and answerability coexist with failed ownership because writing a clinical probe direction moves other findings' answers as much as its own. The same procedure provides handles for natural-image objects, with \cfCocoOwnedPct{}\% ownership, establishing that the deficit is specific to clinical concepts and does not reflect task difficulty. A direction fitted to the model's own answers provides a handle in the same representation. Ownership belongs to the representation and the reader together, with \cfSwapChestOwned{} of \cfSwapChestCells{} chest cells owned across all combinations of tower and reader. Readability does not confer control.
